The supplement provides the analytical expositions, conditional numerical arguments, entire final-cell power outcomes and reproducibility outcomes. Notation is as in the paper: $A,G,c$ are resolved harmful, resolved beneficial and final unresolved masses, respectively, $u$ is final harmful mass and $r=A+G$.

\section{Information decomposition and all boundary cases}
For fixed observable masses, every completion has the joint table
\begin{equation}
 \begin{array}{c|cccc} &1&\cdots&K&\infty\\\hline
 L=1&a_1&\cdots&a_K&u\\
 L=0&b_1&\cdots&b_K&c-u
 \end{array}.
\end{equation}
Its entries are nonnegative and sum to one exactly when $0\le u\le c$. The row harmful marginal is $A+u$. Conditioning on columns gives
\begin{equation}
 H(L\mid J_\Pi)=\sum_{k:a_k+b_k>0}(a_k+b_k)h\left(\frac{a_k}{a_k+b_k}\right)+q(u,c-u).
\end{equation}
This proves the information profile in the article without a selection-model assumption. When $r,c>0$, put $p_R=A/r$, $p_U=u/c$ and $\theta=rp_R+cp_U$. Then
\begin{align}
 \Sprof_\Pi(u)-\tau_\Pi
 &=h(\theta)-r h(p_R)-c h(p_U)\\
 &=r\KL\{\operatorname{Bern}(p_R)\Vert\operatorname{Bern}(\theta)\}
 +c\KL\{\operatorname{Bern}(p_U)\Vert\operatorname{Bern}(\theta)\}\ge0.\label{eq:js}
\end{align}
The equality is interpreted by continuity if a marginal is zero. When $0<r,c<1$, equality holds in the divergence inequality exactly when $p_R=p_U$. From this it follows that $u=Ac/r$. Therefore $\tau_\Pi$ is minimal full information possible and is determined by the resolved timing distribution. Equation~\eqref{eq:js} also presents the model as an information budget that remains after resolving timing for terminal choice.

For interior hidden masses, $h'(x)=\log\{(1-x)/x\}$ and $h''(x)=-1/\{x(1-x)\}$. Hence
\begin{equation}
 \Sprof''(u)=\frac1u+\frac1{c-u}-\frac1{A+u}-\frac1{G+c-u}
 =\frac{A}{u(A+u)}+\frac{G}{(c-u)(G+c-u)}.
\end{equation}
If $r>0$, then at least one numerator is positive, and so the profile is strictly convex whenever $c>0$. If $A=0< G$, then the derivative is $\log\{(G+ c -u)/(c -u)\}>0$ and the minimum is attained at $u=0$. If $G=0< A$, then the derivative is $\log\{u/(A+u)\}<0$ and the minimum is attained at $u=c$. If $c=0$ then the only possible completion is $u=0$ with information $h(A)-T_\Pi=\tau_\Pi$. If $r=0$, $c=1$, then the profile reduces to $h(u)-h(u)=0$ everywhere on $[0,1]$. These cases also explain why minimum-information risk is left undefined from among resolved cases when $r=0$.

For $r>0,c>0$, strict convexity and continuity imply that the sublevel set is empty below its minimum, equal to the singleton of the unique minimizer at equality, and a nonempty compact interval above it. It may touch either boundary; an interior root is not required to exist on both branches. Since every feasible completion is a valid table, and no other completion is possible, its risk interval is sharp. If only coarse observations are given, any completion can be lifted to a hypothetical finer representation with no additional within-coarse-band outcome dependence, so mutual information is unchanged. Fine margins are known to restrict such a lift, explaining why the article phrases sharpness as holding with respect to the supplied partition.

\section{Refinement, strictness and decision frontiers}
Let each finite coarse band $B$ collect fine bands with masses $m_k=a_k+b_k$ and harmful proportions $r_k=a_k/m_k$. For $m_B=\sum_{k\in B}m_k>0$, its harmful proportion is $\bar r_B=\sum_{k\in B}m_kr_k/m_B$. Its entropy gain is
\begin{equation}
 m_Bh(\bar r_B)-\sum_{k\in B}m_kh(r_k)\ge0.
\end{equation}
Strict concavity implies the inequality is strict precisely when the positive-mass portions of $B$ are different. This is summed over coarse bands. The unresolved column stays the same, so this sum is $T_c-T_f$ regardless of its hidden label split. The chain rule for the deterministic map $J_c=g(J_f)$ yields
\begin{equation}
 I(L;J_f)=I(L;J_c)+I(L;J_f\mid J_c),
\end{equation}
and shows that it equals the conditional information. Side-by-side profile comparison yields nesting at all levels of common budget, including empty-set budgets $=\{\}$.

For part 3 of the article's statement, fine compatibility implies $\rho\ge\tau_f>\tau_c$. Hence the coarse right endpoint is strictly right of its argmin. If this endpoint is $<c$, then its coarse profile value is $\rho$ and its fine profile value is $\rho+\Delta$. Fine feasibility excludes that point. Compactness prevents the fine upper endpoint from approaching it without attaining it, establishing the strict inequality. If both endpoints equal $c$, then the upper risk tolerances equalize even though $\Delta>0$. Informative timing alone cannot strictly improve when both bounds are saturated. If $c=0$ then there is no slackening opportunity because there is no hidden risk.

Proof of frontier: 
For the frontier, first note that every compatible set contains $\theta^\dagger=A/r$. Hence $\beta<\theta^\dagger$ intrinsically violates certifiability even if $A\le\beta$; remember $\beta<A$ is the stronger, direct barrier. If $\beta\ge A+c$ then every completion is safe. Otherwise $\theta^\dagger\le\beta<A+c$ which puts $u_\beta=\beta-A$ on the fine right branch of the profile, where it is strictly increasing past the profile minimum. If $\rho\le\Sprof(u_\beta)$ then its sublevel set right endpoint $\le u_\beta$. If $\rho>\Sprof(u_\beta)$ then continuity ensures a feasible point past $u_\beta$ (since $u_\beta<c$), so certification fails. Equality must therefore hold in the certifying direction. Subtracting profiles at $u_\beta$ recovers exact amount of frontier shift. The exception $\beta=\theta^\dagger<A+c$ has certifying budget $\rho=\tau$ only. 

\section{Proof of the simultaneous and projected guarantee}
For a true category probability $p_j\in(0,1)$, write $X_{j,i}=\ind\{Y_i=j\}$. For any fixed $z\in[0,1]$, the process
\begin{equation}
 Z_n(z;p_j)=\prod_{i=1}^n
 \left(\frac z{p_j}\right)^{X_{j,i}}
 \left(\frac{1-z}{1-p_j}\right)^{1-X_{j,i}}
\end{equation}
is a nonnegative martingale, because its next multiplicative factor has conditional expectation $p_j(z/p_j)+(1-p_j)\{(1-z)/(1-p_j)\}=1$. Integrating against a Beta$(1/2,1/2)$ law preserves martingale property by conditional Tonelli, starts at one, and yields the article's $M_{j,n}(p_j)$. At $p_j=0$, almost surely $X_{j,i}=0$ for all $i$, and the limiting mixture is hence nonincreasing with initial value $1$; $p_j=1$ is similar. Hence the maximal inequality gives 
\begin{equation}
 \Pbb\{\exists n:M_{j,n}(p_j)\ge d/\delta\}\le\delta/d.
\end{equation}
The union bound is over categories, not monitoring times; simultaneous containment in every raw interval happens with probability at least $1-\delta$. Thus it suffices to prove the statement on this event. Taking running intersections and intersecting with the simplex both preserve this event. Categories need not be independent within an event. No union bound is taken over monitoring times.
On simultaneous containment, the true summary is in the observable-law region for every $n$, simultaneous containment says that the true terminal harmful mass provides a feasible completion at every $n$. Since its finest-path information cannot be more than $\rho$ and coarse information cannot be negative, this holds because it belongs to the union of the partition-specific risk sets at every $n$. Its path-wise supremum over $n$ is therefore bounded by the aggregate relaxation, and by any conservative numerical enclosure. This proves the time uniform risk guarantee of the article. Since restricting reports to checkpoints or adding evidence gates can only decrease (remove) opportunities to certify, neither require additional error spending. Error cannot be spent to repair feasibility if the sensitivity assumption is violated.
Failure of this feasibility step is not repaired by compatibility of some other completion. The event also implies simultaneous coverage of every element of the partition-specific identified set: on this event the whole true observable-law set is contained in each projection. So the guarantee is strictly stronger than pointwise containment of some chosen terminal completion, while the implemented relaxation may be larger than needed for this guarantee. For several partitions or several streams of risk this same logic would require a simultaneous observable region or an explicitly constructed union allocation; the single partition theorem does not automatically give any such allocation.

\section{Envelope, numerical contracts and branch-and-bound}
For any selected set of categories $B$, simplex membership implies
\begin{equation}
 \max\{\sum_{j\in B}\ell_j,1-\sum_{j\notin B}v_j,0\}
 \le\sum_{j\in B}p_j\le
 \min\{\sum_{j\in B}v_j,1-\sum_{j\notin B}\ell_j,1\}.
\end{equation}
These inequalities give the aggregate mass bounds. For $a,b>0$, $\partial q/\partial a=\log\{(a+b)/a\}\ge0$ and $\partial q/\partial b=\log\{(a+b)/b\}\ge0$. Continuity extends monotonicity to boundaries. Thus evaluating $q$ at simultaneous lower or upper coordinates gives valid lower or upper entropy bounds. The raw upper sum is valid but can be loose; intersecting it with the resolved-mass entropy ceiling and $\log2$ tightens it. A ceiling is not needed to turn the raw sum into an upper bound. This monotonicity argument, rather than attainability of simultaneous upper coordinates, justifies the envelope.

Every categorical summary maps into $\Eset_n$. Since $\Sprof=h(A+u)-T-q(u,c-u)$ decreases in $T$, replacing the actual $T$ by $T_U$ can only enlarge its feasible set. For a search box $B=[A_-,A_+]\times[G_-,G_+]\times[u_-,u_+]$, define
\begin{align}
 c_-&=\max(c_L,1-A_+-G_+,0),&c_+&=\min(c_U,1-A_--G_-,1),\\
 \tilde u_+&=\min(u_+,c_+),&
 z_-&=\max(0,c_--\tilde u_+),\quad z_+=\max(0,c_+-u_-).
\end{align}
Impossible mass constraints discard the box. Otherwise every feasible point has $u\in[u_-,\tilde u_+]$, terminal beneficial mass in $[z_-,z_+]$, and $\theta\in[A_-+u_-,\min(1,A_++\tilde u_+)]$. A lower enclosure is
\begin{equation}
 \underline S_B=\max\left\{0,\ \inf_{x\in[A_-+u_-,\min(1,A_++\tilde u_+)]}h(x)
       -T_U-\sup_{u\in[u_-,\tilde u_+],\,z\in[z_-,z_+]}q(u,z)-\eta_B\right\},
\end{equation}
Here $\eta_B=0$ evaluates in the current source: the desired enclosure is constructed from oriented helper operations instead of an engineered additive buffer. The binary-entropy minimum is at an endpoint of its argument interval. The upper mass-entropy bound is the lesser of an Arb interval upper bound (both arguments are positive), or $(u_++z_+)\log2$. Clamping the information lower bound to zero need not lower-bound the relaxed expression if it is negative, but retains every feasible relaxed point for $\rho\ge0$. Therefore enclosing arithmetic allows $\underline S_B>\rho$ to be a valid rejection criterion.
The numerical implementation has three independent obligations: enclosing elementary mass and entropy arithmetic; preserving likelihood signs through confidence inversion; and retaining every feasible box upon exiting the search. The third property is a corollary of the following search invariant. Arb has an enclosure contract for transcendental operations in its library wrapper (Johansson, 2017; in the article reference list). Confidence inversion and projection in the current accelerated code likewise takes some signs from ordinary binary64 logarithms. The currently available validation leaves that filter without a uniform error analysis. Therefore the complete implemented contract must instead be conditional on these component contracts, rather than proved by the tests themselves.
Inputs are taken to be represented exactly as binary64 dyadics. The rules below assume exact Python integer/rational arithmetic throughout, including correct rounding for conversions to and from binary64 and division by powers of two, correct adjacent-float stepping, and the Arb enclosure contract. Each must perform its numerical function calls sequentially without external modification of its Arb computation context; parallel streams must use separate processes. Confidence inversion and projection check the configured working precision, 128 bits in this campaign. Summary entropy queries use the current Arb context together with outward endpoints, so its enclosure can have a wider enclosure at a different precision setting without altering the mathematical quantity.
\emph{Addition and interval end construction.} To convert an exact rational $v$ to a nearest binary64 dyadic $f$, an endpoint is compared against $v$ exactly and incremented by an adjacent float if necessary to achieve the requested direction. Exact-rational normalization and comparisons are preserved for singleton tables and feasible-box witnesses. Otherwise mass sums, complements, objective sums and gap sums instead evaluate $f=\texttt{fsum}(x_i)$ and a residual $e=\texttt{fsum}(x_i,-f)$. A lower bound is decremented if $e<0$, and an upper bound is incremented if $e>0$. This encloses the exact sum when both sums are correctly rounded, without overflow or failure to account for subnormal residuals. In fact if the exact sum of binary64 dyadics is nonzero then it has magnitude at least as large as the smallest subnormal, so correct rounding of the residual maintains its sign. This strengthens the contract by being explicit, rather than just assuming each floating-point addition is rounded correctly. Endpoints from Arb are yielded as integer--power-of-two pairs; these halves are converted to nearby binary64 dyadics and compared against the exact pair, with outward stepping and a rational fallback in exceptional exponent ranges. Nonfinite quantities produce uninformative endpoints or failure, not merely bit patterns that happen to prefer the desired direction.

\emph{Accelerated confidence inversion.} Let $b=\log B(s+1/2,n-s+1/2)-\log B(1/2,1/2)$ and $t=\log(d/\delta)$. Arb supplies outer binary64 endpoints $o_L,o_U$ for $b-t$. The fast filter evaluates
\begin{align}
 \widehat\ell(x)&=s\,\texttt{log}(x)+(n-s)\,\texttt{log1p}(-x),\\
 e(x)&=2^{10}\epsilon_{\rm mach}
       \{|o_L|+|o_U|-\widehat\ell(x)\}+o_U-o_L.
\end{align}
It declares a positive root sign when $o_L-\widehat\ell(x)>e(x)$ and a negative sign when $o_U-\widehat\ell(x)<-e(x)$. Otherwise it computes the full root expression using Arb; failure to resolve an Arb sign leaves the outer bracket alone. Newton-like root candidates are suggestions. 

The auxiliary verified-bracket routine also has the fast filter, so its name does not suggest that it checks both endpoints against Arb.

The required filter contract is that every accepted sign agrees with the exact mixture-log expression. A complete derivation would need bounds for the platform logarithms, count conversion, products, cancellation and the computed allowance itself. The factor $2^{10}\epsilon_{\rm mach}$ and finite regression checks do not by themselves establish this contract. No incorrect filter sign in a reported experimental stream has been demonstrated here; a uniform implementation guarantee remains unverified.

Conditional on delivering correct signs, convexity provides the usual bracket invariant: A positive sign shifts the outer endpoint, while a nonpositive sign shifts the inner endpoint. If no inner sign is produced, the entire corresponding side is kept. Zero-count and all-count intervals use continuous likelihood values. Running intersections use endpoint maxima/minima. These bracket operations then contain the exact confidence sequence, which is the numerical input hypothesis appearing explicitly in the article's time-uniform theorem.

Summary representation: A category-entropy sum is bracketed both directly and by its simplex complement. The greater lower endpoint and lesser upper endpoint still bracket every law in the intersection of that confidence rectangle with the simplex. Monotonicity of $q$ provides the resolved-entropy bounds, further capped above by the resolved-mass upper bound times $\log2$. Point category tables also use directed aggregate sums and an Arb entropy interval. Their ``singleton'' flag describes the category marginals, not exact real entropy. Approximate normalization accepted by the population data layer is never trusted as an exact simplex identity by the certifier.

Upper-search invariant and exits: Initially, every feasible point of the relaxed projection is in the aggregate box. Directed geometry tests only reject points with impossible mass constraints, and the preceding information lower bound rejects only points above the information budget. A split produces a representable midpoint within its parent interval; its two closed children cover the parent interval, even when rounding causes one child to coincide with its parent. A retained box takes the upward-rounded bound $\min(1,A_++u_+)$. Consequently the current largest active or queued objective upper bound over all boxes encloses the remaining supremum. Incumbents may increase this reported maximum but cannot decrease it. Objective-gap, decision, box-width, stall and node-cap exits preserve the relevant active and queued bounds. Arithmetic failure during a step leaves the active parent retained even if only some children have been added to the queue. Failure before some conservative result is available propagates as a technical failure. When there are no boxes or incumbent remaining, the reported upper bound is one. Thus convergence is not required to dominate $V_n$.

Singleton shortcut: The scalar right-branch search is only for point category inputs with point aggregate mass intervals. It further requires exact equality of those aggregates to the category sums and exact normalization with respect to $c$. Failure of either identity check produces an assumption-free upper bound, without producing an incumbent. Exactly normalized tables with non-point aggregate sums use the box search. Non-normalized point tables go directly to the conservative fallback. For fixed valid $A,G,c$, the larger relaxed endpoint takes exactly the supplied entropy upper endpoint $T_U$, not an independently recomputed nominal entropy. The initial right-branch lower bracket is at or above the exact minimizer $Ac/(A+G)$, and the upper bracket is $c$. Only a certified profile value exceeding $\rho$ may shrink that upper bracket. Convexity on this branch can prove that every feasible value of hidden mass is below it; an ambiguous sign cannot shrink that upper bracket. Upward conversion of $A+u_{\rm hi}$ computes the risk enclosure. Zero resolved or unresolved mass is handled by the continuous boundary formulas. The label \texttt{EXACT\_SINGLETON} is recorded only when a verified lower incumbent equals that upper bound; otherwise the outcome is enclosed or the conservative fallback is applied.

Feasible witnesses and gaps: Candidate allocation and the population root solver only produce proposals. For a proposed table, the certifier exactly sums its provided category entries, defines terminal mass to be the rational complement $c=1-A-G$, checks exact membership of category, aggregate, terminal and hidden-mass in their respective intervals including $0\le u\le c$, demands the actual resolved entropy lie inside its envelope, and requires an Arb upper bound on the actual table's information to be at most $\rho$. These checks produce a normalized version of the original joint table regardless of rounded aggregate metadata provided by the proposal. Its risk is conservatively rounded down before it may become an incumbent. The singleton path also applies these entropy-membership and actual-information conditions after verifying its exact mass identities. Failed or unresolved checks produce no witness. Thus an incumbent cannot exceed the original supremum and also lies below the relaxed supremum. The gap is reported to be rounded upward from enclosure minus this lower incumbent; lack of a witness means lack of any claimed gap. A lower incumbent exceeding $\beta$ refutes certification by this projection, without proving that true risk exceeds~$\beta$.

Rejection searches and final assessment: The compatibility search gives each box a downward-bound equal to $\max\{0,\inf q(A,G)-T_U\}$, which is below every original compatibility floor in that box. The intrinsic search uses a downward enclosure of $A_-/(A_-+G_+)$, which it takes to be zero if its denominator vanishes, and discards boxes only by the same safe geometry and information tests. This ratio therefore bounds $A/(A+G)$ for every compatible original law retained in the box. In both searches the minimum of the active or queued lower bounds is therefore a global lower bound. Approximate point values can terminate either search but can never prune boxes or become the reported lower bound. At any early return condition the active box is included; following splitting its covering children are included. An empty retained queue returns zero rather than a point estimate. Exceptions propagate upward without producing a scientific result. If zero resolved mass is still plausible then the intrinsic bound is not produced. Finally, assessment verifies that the projection and requested sensitivity budget are compatible. Rejections proceed via strict comparisons to upward-rounded budget-plus-guard thresholds, while certification only requires $\widehat U_n\le\beta$ after the evidence gates. Subject to the confidence-sign and summation contracts above, these statements compose correctness of the confidence input, summary enclosure, searches and assessment into correctness of the numerical condition stated as the article's time-uniform theorem.

The current singleton shortcut no longer replaces an approximate population-root midpoint with an upper enclosure. All existing directional regressions continue to support the repaired normalization, entropy and exit paths. They are finite checks, not a proof of the fast sign filter nor an attestation for every saved run.

Strict comparison guards are $10^{-12}$. Population root tolerance is $10^{-12}$; identity residual tolerance is $10^{-10}$; categorical inversion tolerance is $10^{-12}$; nominal outer objective gap is $10^{-6}$ with a 2,000,000-node cap. A normalized box-width exit or stall is reported in addition to objective convergence.

\section{Complete confirmatory outcomes and design}

All rows in Table~\ref{tab:allpower} use 100 independent streams, $\delta=0.05$, and confirmatory indices 1,000,000--1,000,099. Within a law the same finest-resolution streams serve each partition. The ordinary grid is 2,000, 5,000, 10,000, 20,000; the power and matched-control grids and low-prevalence grid are specified in the article. Dense monitoring uses 250, 500, 750, 1,000, 1,250, 1,500, 1,750, 2,000, 2,500, 3,000, 3,500, 4,000, 4,500, 5,000, 5,500, 6,000, 6,500, 7,000, 7,500, 8,000, 8,500, 9,000, 9,500, 10,000. Endpoint-uncertifiable uses 10,000, 20,000, 50,000. Zero counts include all streams and should not be interpreted as indicating zero probability of certification.

\Needspace{6.5in}\small\begin{longtable}{lrrrrrl}\caption{All 32 final-cell certification outcomes. Intervals are conditionwise two-sided exact 95\% intervals for ever certifying by the horizon, using 100 trials per row.}\label{tab:allpower}\\\toprule Scenario & $K$ & $\rho$ & $\beta$ & Horizon & Certified & 95\% interval\\\midrule\endfirsthead\toprule Scenario & $K$ & $\rho$ & $\beta$ & Horizon & Certified & 95\% interval\\\midrule\endhead
Endpoint-uncertifiable & 1 & 0.0800 & 0.070 & 50,000 & 0/100 & [0.000, 0.036]\\
Endpoint-uncertifiable & 2 & 0.0800 & 0.070 & 50,000 & 100/100 & [0.964, 1.000]\\
Endpoint-uncertifiable & 4 & 0.0800 & 0.070 & 50,000 & 100/100 & [0.964, 1.000]\\
Endpoint-uncertifiable & 8 & 0.0800 & 0.070 & 50,000 & 99/100 & [0.946, 1.000]\\
Population decision boundary & 1 & 0.0500 & 0.060 & 20,000 & 0/100 & [0.000, 0.036]\\
Population decision boundary & 8 & 0.0500 & 0.060 & 20,000 & 0/100 & [0.000, 0.036]\\
Ignorable delay & 1 & 0.0200 & 0.090 & 20,000 & 98/100 & [0.930, 0.998]\\
Ignorable delay & 8 & 0.0200 & 0.090 & 20,000 & 0/100 & [0.000, 0.036]\\
Intrinsic impossibility & 1 & 0.0500 & 0.050 & 20,000 & 0/100 & [0.000, 0.036]\\
Low prevalence & 1 & 0.0200 & 0.050 & 20,000 & 100/100 & [0.964, 1.000]\\
Low prevalence & 2 & 0.0200 & 0.050 & 20,000 & 99/100 & [0.946, 1.000]\\
Low prevalence & 4 & 0.0200 & 0.050 & 20,000 & 99/100 & [0.946, 1.000]\\
Low prevalence & 8 & 0.0200 & 0.050 & 20,000 & 99/100 & [0.946, 1.000]\\
No timing control & 1 & 0.0800 & 0.090 & 50,000 & 94/100 & [0.874, 0.978]\\
No timing control & 2 & 0.0800 & 0.090 & 50,000 & 28/100 & [0.195, 0.379]\\
No timing control & 4 & 0.0800 & 0.090 & 50,000 & 0/100 & [0.000, 0.036]\\
No timing control & 8 & 0.0800 & 0.090 & 50,000 & 0/100 & [0.000, 0.036]\\
Misspecified sensitivity & 1 & 0.0200 & 0.045 & 20,000 & 0/100 & [0.000, 0.036]\\
Misspecified sensitivity & 2 & 0.0200 & 0.045 & 20,000 & 0/100 & [0.000, 0.036]\\
Incompatible budget & 4 & 0.0200 & 0.060 & 20,000 & 16/100 & [0.094, 0.247]\\
Incompatible budget & 8 & 0.0200 & 0.060 & 20,000 & 0/100 & [0.000, 0.036]\\
Near degeneracy & 1 & 0.0493 & 0.050 & 20,000 & 99/100 & [0.946, 1.000]\\
Near degeneracy & 8 & 0.0493 & 0.050 & 20,000 & 99/100 & [0.946, 1.000]\\
Dense monitoring & 1 & 0.0200 & 0.065 & 10,000 & 0/100 & [0.000, 0.036]\\
Strong timing & 1 & 0.0800 & 0.090 & 50,000 & 96/100 & [0.901, 0.989]\\
Strong timing & 2 & 0.0800 & 0.090 & 50,000 & 100/100 & [0.964, 1.000]\\
Strong timing & 4 & 0.0800 & 0.090 & 50,000 & 100/100 & [0.964, 1.000]\\
Strong timing & 8 & 0.0800 & 0.090 & 50,000 & 99/100 & [0.946, 1.000]\\
Risk near budget & 1 & 0.0500 & 0.052 & 20,000 & 0/100 & [0.000, 0.036]\\
Risk near budget & 8 & 0.0500 & 0.052 & 20,000 & 0/100 & [0.000, 0.036]\\
Unsafe law & 1 & 0.0500 & 0.040 & 20,000 & 0/100 & [0.000, 0.036]\\
Unsafe law & 8 & 0.0500 & 0.040 & 20,000 & 0/100 & [0.000, 0.036]\\\bottomrule\end{longtable}\normalsize
 

\begin{table}[!htbp]\centering\small
\caption{Current search states among 96 decision-targeted streams at selected checkpoints. C: certified; R: refuted; U: unresolved; S: stopped. Four full-search representatives per cell are excluded. Runtime is the median projection-related checkpoint time over available targeted updates, in seconds.}\label{tab:searchstates}
\begin{tabular}{lrrrrrrr}\toprule
Scenario & $K$ & Events & C & R & U & S & Median (s)\\\midrule
Strong timing, $\beta=.09$ & 1 & 5,000 & 0 & 89 & 7 & 0 & 0.0080\\
Strong timing, $\beta=.09$ & 2 & 5,000 & 96 & 0 & 0 & 0 & 0.6113\\
Strong timing, $\beta=.09$ & 1 & 50,000 & 92 & 0 & 4 & 0 & 4.1088\\
Matched control & 2 & 50,000 & 27 & 1 & 67 & 1 & 13.8887\\
Matched control & 8 & 50,000 & 0 & 94 & 0 & 2 & 0.0981\\
\bottomrule\end{tabular}
\end{table}

Lines add to 96. These are states at the listed checkpoint, unlike headline totals which marginalize over all streams ever certified among the 100, including the four full searches. Decision-directed halts occur when some bound establishes one side of the risk-budget comparison; full searches continue toward their objective-gap criterion. Full search may report a stall even after having established a certifying upper bound. Here, refute counts decision-refuted terminations where there was a numerically verified feasible lower incumbent above $\beta$; it does not certify nor does it imply that true risk exceeds~$\beta$. The subtlety depends on the numerical conditions listed above; termination labels by themselves do not certify them. Unresolved counts stalls where search neither certified nor provided a refutation; subsequent search may certify or may not. Stopped has no current update after its running confidence region depleted, but still counts against denominator events. None of the runs posted a model-incompatible state at any of the 160 power checkpoints.
{\sloppy The supplementary laws use parameter ordering $(\theta,q_1,q_0,\lambda_1,\lambda_0)$. High unresolvedness uses $(.05,.70,.40,.35,-.10)$. Timing-only uses $(.05,.10,.10,.45,-.15)$ and terminal-only uses $(.05,.30,.05,0,0)$. Near degeneracy uses $(.01,.90,.01,.80,-.80)$; intrinsic impossibility uses $(.15,.10,.10,0,0)$. Unsafe, risk-near-budget, incompatible-budget, misspecified-sensitivity and population-decision-boundary use the standard timing-and-terminal harmful-late law from the main text. Ignorable-delay and dense-monitoring use independent-resolution. 
The 32-band validity law was refreshed at generation at 32 bands using the standard timing-and-terminal parameters, rather than being a coarsening of the eight-band law.\par}



In all seven validity scenarios $\rho=I_{\rm true}+0.01$, except minimum-information completion, which sets terminal harmful mass equal to $Ac/r$ and takes $\rho=\tau+0.002$. The near-certification risk budget is itself a sharp upper bound at the listed $\rho$ plus 0.002; other validity budgets are listed as 0.05. Each scenario simulated 500 independent streams, seed indices 0 through 499, for up to 500 matured events each. Nominal error budget is unchanged. Construction exclusions are monitored at every event and five risk-bound diagnostics are available at the five 100-event checkpoints.

A separate descriptive assessment of the longer-horizon power records retains the longest observed horizon for each of 21 distinct law--partition stream sets. Across 2,100 indexed stream outcomes, 58 exclusions of the true categorical law were recorded (2.8\%), with at most six exclusions among 100 streams in any set. These observations are not 2,100 independent Bernoulli trials: some partitions share event streams, and empty confidence regions terminate further updates. Accordingly, this aggregate is descriptive and is not assigned a pooled binomial confidence interval or a campaign-wide coverage interpretation.

For count $x$ from $m$ streams, lower endpoint of the two-sided 95\% binomial interval is zero if $x=0$ and otherwise $B^{-1}_{x, m-x+1}(0.025)$, and upper endpoint is one if $x=m$ and otherwise $B^{-1}_{x+1, m-x}(0.975)$. One-sided violation boundaries equal one if $x=m$ and otherwise use $B^{-1}_{x+1, m-x}(0.95)$. These are Monte Carlo confidence intervals distinct from the stream's own time-uniform confidence region. Prespecified final-horizon power target plugs into this interval, taking the two-sided lower endpoint, just as the stored publication artifacts do.

Repeated-static monitoring produces per-category Wilson intervals with $z=\Phi^{-1}\{1-\delta/(2d)\}$, and no multiple looks correction. It thus can attain greater power without sharing TrajCert's uniform-time guarantee. Applied to the final horizon of valid ignorable-delay, its counts are 98/100 at $K=1$ and 0/100 at $K=8$ while the ignorable-delay benchmark certifies 98/100 for either partition. In unsafe law, the latter benchmark falsely certifies for 7/100 samples at each displayed partition; its independence assumption is violated there. These Monte Carlo diagnostics flag assumption violations. Dense-monitoring sees that benchmark certify for 97/100 streams whereas neither sensitivity projection certifies; the stronger identifying assumption changes the population target bound.
TrajCert power labels record zero false risk certificates. Stored premature-certification counter is zero: it only counts certificates from cells marked as compatible but not population certifiable, so does not include model-incompatible allocations. Sensitivity-misspecified two-band cell is population certifiable under the assumed model but has larger true risk $0.05>0.045$ and larger true partition information than $\rho=0.02$; observed zero certifications do not prove that model correct. Incompatible four-band cell certifies 16/100 despite $\rho<\tau$ in that population, for risk below its budget. Both illustrate gap between risk error and certificate issued under an incorrect structural hypothesis.
\section{Secondary numerical and transport diagnostics} 
Node-budget sensitivity run at same finite-sample level produces outer bound $0.356365$ on all seven caps from 1,000 to 2,000,000 nodes. Every single search terminates by box at 104 nodes with gap $0.157533$. Increasing budget cannot help search stopped by other rule. This finite-sample phenomenon documents a limitation of the relaxation/stopping rules pairing, not evidence of optimizer convergence or numerical stability.
Foreign information analysis evaluates local risk decisions against altered population risk tables. Path substitution maintains local endpoint masses while grafting conditional timing masses from another law; pooling takes simple average of local and foreign law masses. Heterogeneity linearly interpolates the five law parameters from local to foreign values. Table~\ref{tab:foreign} shows every effect category over 280 design cells. At heterogeneity 100\%, pooling induces relative error to local analysis in 63 cells, substitution in two. These are deterministic grid-design sums, not Monte Carlo estimates of error probability.

\begin{table}[!htbp]\centering\small
\caption{Foreign-information effects relative to the local analysis. Each row sums to 280. NC: no change; BT: tightening without a newly introduced error; CW: widening; EB: newly introduced erroneous bound; MI: model incompatible. Labels classify relative effects, not validity under the local true law.}\label{tab:foreign}
\begin{tabular}{rlrrrrr}\toprule Heterogeneity & Replacement & NC & BT & CW & EB & MI\\\midrule
0.00 & Path & 271 & 0 & 1 & 0 & 8\\
0.00 & Pooled & 280 & 0 & 0 & 0 & 0\\
0.25 & Path & 102 & 61 & 87 & 1 & 29\\
0.25 & Pooled & 31 & 80 & 110 & 28 & 31\\
0.50 & Path & 102 & 53 & 87 & 3 & 35\\
0.50 & Pooled & 31 & 68 & 98 & 48 & 35\\
0.75 & Path & 102 & 43 & 87 & 3 & 45\\
0.75 & Pooled & 31 & 62 & 91 & 60 & 36\\
1.00 & Path & 114 & 42 & 72 & 2 & 50\\
1.00 & Pooled & 30 & 57 & 94 & 63 & 36\\
\bottomrule\end{tabular}

\end{table}

An introduced error is the loss of true-risk containment, an upper bound below the true risk, or a false population certificate when that error was absent locally. This grid further contains budgets that fail even for the local law. For instance with full heterogeneity there are 62 and 106 containment failures in the path and pooled tables respectively, but only two and 63 introduced errors respectively. Therefore a ``tightening'' or ``widening'' label does not confirm validity. With zero heterogeneity, naive pooling leaves all 280 cells unchanged; path replacement regenerates shapes at analysis resolution and may therefore differ from simple coarsening of the finest local law. It is therefore not equivalent to a pure zero-change control. These diagnostics motivate using trajectories from the target population, but do not by themselves constitute a transportability theorem. Full outcome counts and both underestimation and false-certification indicators are archived with the evidence record.

\section{Reproducibility record} 
The seeds, generated from labeled streams, are documented at time of simulation. The design and analysis documented here have been successfully executed with these seeds. The declared environment successfully reproduced these batches using the same seeds.

The documented design and reported results used the same evaluation configuration and dependency versions. Specifically, they used the same laws, budgets, grids, seed indices and power/threshold pairs. The batch size increased from 50 to one for coverage. The power batch size increased from ten to one. The completed reproducibility record contains 3,200 power batches and 3,500 coverage batches; batches do not span streams and each contains one stream. Although reproducibility is tied to dependency fingerprints, the ones recorded omit their source tree. The dependency fingerprint for an installed tree can be independently matched to historical streams if its predecessor hierarchy is also recorded, but this record does not contain enough information to do that. Selected-case equivalence and feasible oracle metrics are local to the reported batches so do not imply equivalence of the full evaluation.

The declared environment is Python 3.12 or later, NumPy 2.5.2, SciPy 1.17.1, mpmath 1.3.0 and python-flint 0.9.0. Dependency locks are distributed with the code. System information on the physical CPU or current machine load is not included with the timing records.

Streams were seeded by truncating the SHA-256 hash of a semantic event stream namespace and stream index to integer initial states for NumPy's PCG64 random generator. Namespace bits are intended to protect against several categories of cross-stream equivalence. The namespace includes the law and finest-band count but not analysis partition, sensitivity budget, risk tolerance budget or experiment name. Stream indices are equal for coarse analyses of the same budget sharing the same prefix of the finest. Optional stopping in the independent-panel grid used twenty streams rather than ten, by a posted design amendment. Pilot work with ten streams selected events from the ordinary population decision example but found no certification through 200,000 matured events. It is not included in the 100-stream confirmatory batches. Failures at earlier horizons do not by themselves establish that certification is impossible at later horizons. Pilot batches whose results informed this work are private. Their accompanying design record is not an external preregistration.

The reproduction evidence supplement contains the scientific reproducibility record, batches of individual power and validity estimates, numerical tables used to generate figures in the paper, figure source data, configuration and dependency locks, with checksums. External raw diagnostic data and pilot batch records are not distributed. Successfully reproducing paper assets from this evidence record does not require an additional simulation.
